\section{Discussion and limitations}
\label{sec:discussion}

The method computes the present-time field of a uniformly moving, TSC-shaped
bunch as a free-space FFT convolution. It is a far-field method.

\paragraph{What it is good for.} In the far field the relative error against
the point field is below $10^{-3}$ on the beam axis from $d=6$ for every
$\gamma$ \citep{N2b}. Against the TSC-convolved continuum: $\gamma=1$, aspect
1 crosses $3\times10^{-4}$ at $d=12$ ($1.1\times10^{-4}$); $\gamma=1$, aspect
4 at $d=32$ ($2.6\times10^{-4}$); $\gamma=10$ is $6.5\times10^{-3}$ at $d=32$,
not below $3\times10^{-4}$ within $48$ cells \citep{W4b,V2b}. At
$2\times10^{6}$ particles in a $20$-cell box the compute-only total is
$64.8$\,ms, dominated by the deposit ($56.5$\,ms), and flat in $\gamma$
\citep{SB1}. Repeated runs are bit-identical (md5 \texttt{d3641d64\ldots});
the Coulomb limit reproduces the TSC-convolved Coulomb field to
$4.88\times10^{-14}$ \citep{W4b,coulomb}.

\paragraph{What it is not.} The mesh alone does not reach $10^{-3}$ where the
slices cut the bunch core: on the $N=1000$ slice geometry the mesh-only worst
per-component plane error is $0.12$--$0.41$ for $\gamma=1$ and
$0.11$--$4.60$ for $\gamma=10$ (Table~\ref{tab:slices}). The exact near shell
(Section~\ref{sec:nearsheet}) is opt-in and its assembly is exact to roundoff;
at the measured $R=3/4$ radii on the $h=1$ slice grid the residual is the mesh
outside the shell, not the tile subtraction
(Section~\ref{sec:slices-runs}). The perpendicular direction at $\gamma=10$
does not converge at the natural cell size ($1.5\times10^{-2}$ at $d=48$,
aspect 1; the TSC smoothing alone is $1.48\times10^{-2}$,
Table~\ref{tab:d48}), and reaching $10^{-3}$ perpendicular at $r=8$ cells would
need $h\approx1/29$, a $2.4\times10^{4}$-fold increase in $M^{3}$ FFT work
\citep{N2b,D11}. At aspect 4, $\gamma=1$ already fails the gate at $d=48$
($2.4\times10^{-3}$) \citep{N2b}. The cross-validation confirms the same
boundary: the mesh does not agree with the exact reference at the reference
floor; the disagreement is per-source, in the TSC support, and does not grow
with particle count \citep{W7t,D13}.

\paragraph{Performance.} The cost is two regimes with a crossover set by the
box, not the particle count \citep{SB1}. At large $N$ the deposit dominates,
linear in $N$, $52$--$57$\,ms on the $2\times10^{6}$ rows. At large boxes the
FFT and kernel dominate: from $M=60$ to $M=220$ the total goes from
$64.8$\,ms to $233$\,ms, and halving $h$ costs $1.07\times$ and $2.35\times$
\citep{SB1}. The compute timings are single-GPU on one \texttt{rosi\_v100}
host. The output path has a multi-rank measurement: at $512^{3}$, tolerance
$10^{-3}$, the chunked writer runs in $54.3$\,s for one rank and $43.2$\,s
for two, the six datasets bitwise identical
(Section~\ref{sec:mpi-scale}) \citep{FIX1}. The near-field cost is a separate
wall: per-pair, $5.4\times10^{-5}$\,s for the ported stationary box8/rem8
evaluator, $3.1\times10^{-4}$\,s out of support and
$6.9\times10^{-3}$\,s in support for the accurate moving evaluator,
$1.62$\,s for the upstream adaptive settings
(Section~\ref{sec:cost}) \citep{K3}.

\paragraph{Scope.} The defensible scope, in floored relative per-component
error against the continuum, is $\gamma=1$ aspect 1 from $d=12$, $\gamma=1$
aspect 4 from $d=32$, and the axial component at $\gamma\le10$
\citep{W4b,N2b,D11}. The blanket statement ``$\gamma\le10$ meets $10^{-3}$''
is not supported: it fails perpendicular at the natural cell size,
$6.5\times10^{-3}$ at $\gamma=10$, $d=32$, aspect 1, still $1.8\times10^{-3}$
at $d=48$, not below $3\times10^{-4}$ within $48$ cells
\citep{W4b,N2b,V2b}. Each statement in Section~\ref{sec:accuracy} names which
metric (floored, unfloored, absolute) supports it; the floored column is
electric-only at $\gamma=10$ because the floor excludes the magnetic field,
and the judge metric $|\mathrm{ref}-\mathrm{mesh}|/(1+\max)$ reports $d=2$
for every case and does not discriminate. The radii in
Table~\ref{tab:crossover} are axis directions only; oblique directions are
$\sim1.5\times$ larger. Scale invariance
(Section~\ref{sec:scaleinvariance}) means only $r/h$ moves the error, so
reaching $10^{-3}$ at $r=8$ needs $h\approx1/29$
($2.4\times10^{4}$ times the $M^{3}$ nodes) \citep{N2b,D10}.
